\chapter{Known limitations in version 1.23.0}\label{ch:knownissues}

No program does everything perfectly. This chapter lists the limitations of
ProPrep version 1.23.0 that we met while writing this manual, so that they do
not surprise you. Each one comes with a way around it, where we found one.
Everything here was seen on a real computer unless the text says otherwise.

The tables follow the order of the main menu. For help with a particular
message on the screen, see Chapter~\ref{ch:troubleshooting}.

\section{Software that ProPrep does not include}\label{sec:knownissues-software}

Some steps need programs or services that are not part of ProPrep. Without
them, ProPrep gets as far as it can and then stops or skips that part.
Appendix~\ref{app:software} explains each one in full.

\begin{longtable}{@{}>{\raggedright\arraybackslash}p{0.20\linewidth}>{\raggedright\arraybackslash}p{0.40\linewidth}>{\raggedright\arraybackslash}p{0.34\linewidth}@{}}
\toprule
\textbf{Missing} & \textbf{What happens} & \textbf{Workaround} \\
\midrule
\endhead
\bottomrule
\endlastfoot
MODELLER license key &
  No loops, missing residues or atoms built (\key{4c}); standard mutations not
  applied (\key{4b}); no alternate-location choice; no homology models
  (\key{3a}). &
  Free academic key (Chapter~\ref{ch:installation}). tLEaP still adds missing
  side-chain atoms. \\
\texttt{pmemd} (Amber), GPU &
  The default MD engine is missing; runs stop at once. &
  Choose \texttt{sander} or \texttt{sander.MPI}, which are included
  (Chapter~\ref{ch:runningmd}). \\
Gaussian &
  No RESP charges for ligands or modified residues (\key{2c}); Gaussian QM/MM
  input is written but cannot be run. &
  AM1-BCC charges made outside ProPrep can be imported
  (Chapter~\ref{ch:forcefields}). \\
ORCA &
  The ORCA route of \key{5c} writes no files, because it needs ORCA's
  \texttt{orca\_mm} program. &
  Install ORCA, or use the Gaussian route (Chapter~\ref{ch:qmmm}). \\
Internet connection &
  PDB and AlphaFold downloads, BLAST, AlphaFill and the OPM comparison fail;
  some failures are reported misleadingly. &
  Work online for these steps. \\
\end{longtable}

\section{Start-up, menus and the project folder}\label{sec:knownissues-start}

\begin{longtable}{@{}>{\raggedright\arraybackslash}p{0.16\linewidth}>{\raggedright\arraybackslash}p{0.44\linewidth}>{\raggedright\arraybackslash}p{0.34\linewidth}@{}}
\toprule
\textbf{Tool} & \textbf{What happens} & \textbf{Workaround} \\
\midrule
\endhead
\bottomrule
\endlastfoot
Start-up &
  The \prompt{MODELLER not found ... To install} message appears whenever the
  license key is missing, even though MODELLER is installed. &
  Add the key, or ignore the message (Chapter~\ref{ch:installation}). \\
 &
  An \texttt{AMBERHOME} already set in the Terminal is kept without a warning,
  so another Amber installation may be used. &
  Check \texttt{echo \$AMBERHOME} before starting
  (Chapter~\ref{ch:installation}). \\
 &
  A project name containing \texttt{/} stops ProPrep. \typed{undo} typed at
  the project-name question creates a folder called \file{undo}. &
  Use simple names (letters, digits, \texttt{-}, \texttt{\_}). \\
\midrule
\texttt{proprep-web} &
  Closing the last browser tab shuts ProPrep down after about two seconds,
  without asking. &
  Exit with \key{x} and download the project first. Files and the session
  are kept either way. \\
 &
  \textsf{Download project} packs everything in the folder
  \texttt{proprep-web} was started from. &
  Start \texttt{proprep-web} from a dedicated work folder, not your home
  folder. \\
\midrule
All tools &
  Each tool picks its input file by a fixed order of preference, not by date.
  Files ProPrep made come before a newer download. &
  Read the \prompt{Using ... file:} line. One protein per project folder
  (Chapter~\ref{ch:projectfolder}). \\
 &
  Fixed output names such as \file{filtered_structure.pdb} and
  \file{structure_with_md_names.pdb} are overwritten without asking. &
  Copy or rename versions you want to keep. \\
 &
  Some hints have a gap where a key should be, such as \prompt{or press\quad to
  return to the main menu}. &
  The missing key is \key{m}. \\
 &
  Some steps report success although the result is missing or empty. &
  Check the files in the project folder. \\
\midrule
Main menu &
  A \texttt{✓} means ``available'', not ``done''. Codes are case-sensitive in
  the Full Menu. &
  Look for \texttt{●} inside the tool's own menu. Type codes in lower case. \\
 &
  Answering \key{n} to \prompt{Are you sure you want to exit?} inside a tool
  returns to the main menu. &
  Open the tool again. \\
\midrule
Workflow Mode &
  Stage marks in the guide bar are not always right, and its help says
  \key{s} skips a stage (it jumps to Simulate). &
  Files are not affected (Chapter~\ref{ch:menus}). \\
\midrule
\key{w} Workspace &
  \textsf{Save workspace} stores almost nothing; \textsf{Load workspace} and
  \textsf{Load command history} restore little. \textsf{Reset workspace} also
  leaves Workflow Mode. &
  Rely on session files and the project folder. Press \key{g} to return to
  Workflow Mode. \\
\midrule
\key{h} Help &
  The Help page lists the menu codes of an older version. &
  Use Appendix~\ref{app:menumap}. \\
\midrule
\key{f} Feedback &
  The default answer to \prompt{Submit this feedback to GitHub Issues?} is
  \key{y}. &
  Type \key{n} if you are not ready. \\
 &
  The form's opening box points to option~2 of Preferences; it is option~3.
  The token is saved as plain text in \file{~/.proprep/settings.json}. Its list
  of tools uses old names. &
  Use \key{p} then \key{3}. Give the token minimal permissions. Choose
  \textsf{0~General} and name the tool (Chapter~\ref{ch:help}). \\
\end{longtable}

\section{Loading structures (1)}\label{sec:knownissues-load}

\begin{longtable}{@{}>{\raggedright\arraybackslash}p{0.16\linewidth}>{\raggedright\arraybackslash}p{0.44\linewidth}>{\raggedright\arraybackslash}p{0.34\linewidth}@{}}
\toprule
\textbf{Tool} & \textbf{What happens} & \textbf{Workaround} \\
\midrule
\endhead
\bottomrule
\endlastfoot
\key{1a} Structure Loader &
  Entries available only in mmCIF format cannot be loaded. File browsers
  accept only \file{.pdb} files, so AlphaFold and AlphaFill \file{.cif}
  files are not listed. &
  Load AlphaFold models by UniProt accession, or download the PDB-format file
  (Chapter~\ref{ch:structure}). \\
 &
  Without internet, AlphaFold says \prompt{No AlphaFold structure found}. &
  Check the connection. \\
 &
  AlphaFold metadata gives the number of atoms as \prompt{Length: ...
  residues}. After a gene-name search, the UniProt ID shows as \prompt{N/A}. &
  Read pLDDT, not the length. Retrieve by UniProt accession. \\
 &
  AlphaFill says \prompt{No transplanted ligands/cofactors found} even when the
  model has them. &
  Look at the file in the viewer (\key{ua}). \\
 &
  Downloading several search results at once skips the summary and the
  hydrogen question, and only the last entry is used. &
  Download each entry on its own. \\
 &
  The hydrogen-removal answer is reused for the rest of the session. Some
  hydrogens of alternate locations are not removed (3NIR: 18 left). &
  Start a new project to answer differently. \\
\midrule
\key{1b} Biological Assembly Generator &
  \key{m} (manual matrix) stops ProPrep with \prompt{name 'self' is not
  defined}. &
  Choose a listed assembly number instead. \\
 &
  After building, \key{1b} uses the assembly as input and offers it again.
  On a file without a header it finds \prompt{No REMARK 350 records}. &
  Run \key{1b} once, right after loading; cancel with \key{c}. \\
 &
  An assembly downloaded directly (structure type~2) is a two-model file with
  repeated chain names. &
  Prefer \key{1b} for simulations. \\
\end{longtable}

\section{Detecting sites and making parameters (2)}\label{sec:knownissues-detect}

\begin{longtable}{@{}>{\raggedright\arraybackslash}p{0.16\linewidth}>{\raggedright\arraybackslash}p{0.44\linewidth}>{\raggedright\arraybackslash}p{0.34\linewidth}@{}}
\toprule
\textbf{Tool} & \textbf{What happens} & \textbf{Workaround} \\
\midrule
\endhead
\bottomrule
\endlastfoot
\key{2a} Redox Site Detector &
  Caps and crystallization additives (ACE, GOL, ACT, EDO) are listed as
  organic cofactors. &
  Select only the centers you want. \\
 &
  Atom numbers in the bond search refer to a list that is sorted again for
  each atom; bond lengths are not checked (an 8.21\,Å ``bond'' was
  accepted). &
  Type atom names such as \typed{SG} instead of numbers. \\
 &
  The single-site route does not show the bonds a heme recipe needs, so
  \key{4d} may find no compatible transformer. &
  In \key{4d} type \typed{why 1}, then add the bonds; or use templates
  (Chapter~\ref{ch:redox}). \\
 &
  In the site editor, a non-number at a residue-number prompt closes the
  editor and loses the edits. &
  One answer per prompt; type \typed{save} often. \\
 &
  Export rejects the whole list if one number is invalid. Import also lists
  session files, which contain no sites. &
  Use numbers 1--3. Import \file{<PDB>_redox_sites.json}. \\
\midrule
\key{2b} Force Field Explorer &
  The default selection is ff14SB with TIP3P, although the same screen
  recommends ff19SB with OPC. Invalid numbers are ignored silently. &
  Type the numbers you want; check the \prompt{Loading:} line. \\
 &
  The atom-type list starts with numbers such as $-0.00006$. Typing \key{q} at
  a filter searches for ``q''. &
  Always type a filter. Use \key{b} to go back. \\
 &
  Your own library entries appear under \textsf{Redox Cofactors (built-in)},
  which shifts the row numbers. The custom file loader cannot take a typed
  path. &
  Choose rows by name. Copy files into the project folder first. \\
\midrule
\key{2c} Force Field Parameterizer &
  Every charge route needs Gaussian; there is no AM1-BCC option. Without the
  Gaussian output, step~4 prints a technical report (ProPrep has not
  crashed). &
  Run Gaussian and resume, or import AM1-BCC parameters made outside ProPrep
  (Chapter~\ref{ch:forcefields}). \\
 &
  Option~4 may say \prompt{No parameterizations found in workspace} although
  a job was saved. &
  Use option~3 and resume. \\
 &
  Several copies of one ligand share one folder; answering \key{n} for the
  second copy deletes the first copy's work. &
  Parameterize one copy only. \\
 &
  For a modified amino acid, the Gaussian inputs start from the parent residue
  (SEP is written as plain serine). CSD is mapped to ALA. &
  Check the Force Field Explorer first (SEP is already covered). Edit the
  \file{.gjf} files and the mapping. \\
 &
  \textsf{Save as a new version} in the library makes the new version the
  default and reuses the same file names. &
  Rename the new files before importing a second version. \\
 &
  The built-in help calls the metal steps ``coming soon''. &
  The 15-step metal checklist is current. \\
\midrule
\key{2c} metal sites &
  Step~12 can find no residues and write empty Gaussian inputs. No
  workaround was found. &
  Use a ready-made set (zinc, hemes, iron--sulfur) where possible. \\
 &
  Step~10 fails with \prompt{Fatal Error! bond: Argument \#2 is of type
  String} when \key{2a} sites with bonds exist. &
  Do not run \key{2a} beforehand; let step~11 detect the site. \\
 &
  Chloride ions are classed as unknown; supplying an ion library renames them
  to silver. &
  Discard crystallization ions in step~1. \\
 &
  Step~12 asks for a formal charge for every ligand atom. The suggested model
  charge is always 0. Protonation suggestions ignore metal binding and
  disulfides. &
  Answer 0 (\,$-1$ for a thiolate S). Work out the total yourself. Use HID,
  CYX or CYM where needed. \\
 &
  A metal on a crystal symmetry axis gives \prompt{No nearby atoms found}. &
  Choose a fully coordinated metal, or build the assembly first. \\
\midrule
Library command &
  With the conda installer, \texttt{proprep-library} stops with \prompt{bad
  interpreter}. &
  Type \texttt{python -m} \texttt{proprep.utils.\allowbreak
  user\_library\_\allowbreak snapshots} with the same options. \\
\end{longtable}

\section{Comparing structures (3)}\label{sec:knownissues-compare}

\begin{longtable}{@{}>{\raggedright\arraybackslash}p{0.16\linewidth}>{\raggedright\arraybackslash}p{0.44\linewidth}>{\raggedright\arraybackslash}p{0.34\linewidth}@{}}
\toprule
\textbf{Tool} & \textbf{What happens} & \textbf{Workaround} \\
\midrule
\endhead
\bottomrule
\endlastfoot
\key{3a} Homology Searcher &
  A BLAST search has no time limit and cannot be cancelled from the menu;
  ours took almost four hours. &
  Search \texttt{swissprot} or \texttt{pdb} with few hits. Press
  \key{Control}+\key{C} to stop. \\
 &
  Homology models (option~5) need MODELLER. AlphaFold models of hits need a
  search of \texttt{nr} or \texttt{swissprot}. &
  See Section~\ref{sec:knownissues-software}. \\
\midrule
\key{3b} Structure Aligner &
  After a TM-align run, viewing the results fails with an \prompt{Error
  executing ...} line. &
  Open the aligned files in the viewer. \\
 &
  The file browser mentions CIF but accepts only \file{.pdb}. Loaded structures
  are not offered after you decline earlier aligned files. &
  Add structures with the file browser. \\
 &
  Option~2 asks you to run redox detection in the ``PDB Filter module''; the
  tool meant is \key{2a}. A wrong residue range is skipped, not asked again. &
  Run \key{2a} first. Run the alignment again. \\
 &
  An aligned target you select is used by later tools in preference to newer
  downloads. &
  Answer \key{n} at target selection if you only want to compare. \\
\end{longtable}

\section{Preparing the structure (4)}\label{sec:knownissues-prepare}

\begin{longtable}{@{}>{\raggedright\arraybackslash}p{0.16\linewidth}>{\raggedright\arraybackslash}p{0.44\linewidth}>{\raggedright\arraybackslash}p{0.34\linewidth}@{}}
\toprule
\textbf{Tool} & \textbf{What happens} & \textbf{Workaround} \\
\midrule
\endhead
\bottomrule
\endlastfoot
\key{4a} PDB Filter &
  A range (\typed{1-50}) or an out-of-range number at \prompt{Select residues}
  stops the filter, and all answers are lost. &
  Type single numbers separated by commas, or \typed{all}. \\
 &
  The water question accepts one choice only. In a water range, every number
  must be a water. &
  Type one choice; list water numbers. \\
 &
  Answering \key{n} at the final confirmation returns without a message.
  Alternate locations are not removed. &
  Run the filter again. Alternate locations need \key{4c} with MODELLER. \\
\midrule
\key{4b} Amino Acid Mutator &
  Standard mutations are not applied without MODELLER, and no message says
  so. \typed{clear} deletes all pending mutations without asking. &
  Get a MODELLER key. Type carefully. \\
\midrule
\key{4c} Structure Fixer &
  Without MODELLER only caps are added, and nothing is written unless at least
  one cap is chosen; a pending non-standard mutation is then dropped. &
  Add at least one cap (Chapter~\ref{ch:cleaning}). \\
 &
  The repaired file is numbered from 1 (2RH1: 29--1161 becomes 1--506). &
  Use the new numbers in later steps. \\
 &
  The template comparison reports a missing OXT on nearly every residue and
  takes 1--2 minutes. The sequence check counts numbering jumps in
  engineered proteins as missing residues, or may report no problems when it
  skipped its checks. &
  Use REMARK~465 and REMARK~470 in the PDB file. \\
\midrule
\key{4d} Redox Site Preparer &
  Microstates can be generated only once per session. &
  Start a new session in a new project folder. \\
 &
  Each heme recipe needs one heme per site (4HHB grouped four). &
  Detect one heme per site. \\
 &
  Progress lines and the residue mapping use intermediate residue numbers.
  New propionate residues can be numbered between crystal waters. &
  Trust \file{final_transformed_structure.pdb}. Check the QM/MM residue list
  by name (\key{5c}). \\
\midrule
\key{4e} Protonation State Analyzer &
  An invalid pH, threshold or override is replaced by the default without
  asking again. &
  Run the option again. \\
 &
  After overrides, \textsf{View} and \textsf{Export} still show PROPKA's
  states and net charge. The export prompt does not list its formats. &
  The written PDB file is correct. Type \typed{text} or \typed{csv}. \\
 &
  Disulfide cysteines are not renamed to CYX. Oriented files are not used. &
  Use \key{2a} and \key{4d} for CYX. Orient last. \\
\midrule
\key{4f} Structure Orientator &
  An oriented file is always the input, so rotations add up and names grow
  (\file{..._oriented_oriented_...}). &
  Orient once, as the last preparation step. \\
\midrule
\key{4g} Membrane Builder &
  Settings return to their defaults, without a message, after choosing another
  protein, filtering, or a build. &
  Review (\key{r}) before every build. \\
 &
  If tLEaP cannot add hydrogens, ProPrep offers packmol-memgen's own
  protonation; in our test that build was refused minutes later. &
  Remove ions, waters and ligands with \key{4a} first. \\
 &
  A solvate-only build is still called \file{bilayer.pdb} and lists a lipid
  force field. &
  Ignore; the system has no lipids. \\
\end{longtable}

\section{Simulation and analysis (5)}\label{sec:knownissues-simulate}

\begin{longtable}{@{}>{\raggedright\arraybackslash}p{0.16\linewidth}>{\raggedright\arraybackslash}p{0.44\linewidth}>{\raggedright\arraybackslash}p{0.34\linewidth}@{}}
\toprule
\textbf{Tool} & \textbf{What happens} & \textbf{Workaround} \\
\midrule
\endhead
\bottomrule
\endlastfoot
\key{5a} Topology Generator &
  The default answer \key{n} to \prompt{continue without bond definitions?}
  discards the input. &
  Type \key{y} for an ordinary protein. \\
 &
  Success messages appear even when tLEaP fails, and an empty
  \file{.prmtop} is counted as made. A membrane build counts an internal file
  as a topology. &
  Trust \prompt{Errors =}, the file size and \file{leap.log}. \\
 &
  Two inputs write the same output names; a failed build empties the earlier
  topology. &
  Give each version its own names (\texttt{saveamberparm} line). \\
 &
  There is no warning before building with an unprepared cofactor, with
  disulfides still named CYS, or with CYX but no bond lines. &
  Prepare cofactors first; use CYX \emph{and} bond lines
  (Chapter~\ref{ch:topology}). \\
 &
  Implicit-solvent builds get no \texttt{set default PBradii} line and fail
  with crystal waters. &
  Add the line in the template; remove waters with \key{4a}. \\
 &
  No warning for ff19SB with TIP3P. DNA, RNA, sugar and small-molecule force
  fields are loaded by default. &
  Choose OPC with ff19SB. Type \key{1} (None) for force fields you do not
  need. \\
 &
  The PDB file is edited in place to add TER lines. &
  The original is kept as \file{.pdb.ter_backup}. \\
 &
  Option~3 (all microstates) can end with no files and no message. After
  \key{4d} batch generation, option~2 is blocked. &
  Run option~2 on each microstate file, in a separate project. \\
 &
  A constant-pH \file{.cpin} includes ASP, GLU and HIS only with the
  constant-pH force field and renamed residues. Afterwards the cpin topology
  becomes the default for \key{5b}. &
  Build as described in Chapter~\ref{ch:topology}. Answer \key{n} to
  constant-pH MD for normal runs. \\
 &
  The bond editor does not check what you type, and \key{q} saves. A loaded
  topology pair cannot be used by options 2--6. &
  Check the bond list before saving. \\
\midrule
\key{5b} MD Manager &
  Every stage starts with \texttt{pmemd}, which is not installed; the
  recommended assignment picks \texttt{pmemd.cuda} even without a GPU. &
  Choose \texttt{sander} at the hardware step (Chapter~\ref{ch:runningmd}). \\
 &
  The review table keeps the original stage lengths. Edits to settings are not
  checked. &
  Check \file{simulation.mdin} in each stage folder. \\
 &
  \textsf{Load a saved protocol} always says \prompt{No saved protocols
  found}. &
  Import the protocol's \file{.json} file. \\
 &
  A minimization made with the wizard includes pressure-control settings. &
  Use a predefined stage, or answer the pressure questions carefully. \\
 &
  Distance restraints: the import reports an error although it worked; typed
  entry fails because the file has no chain names; only \file{.disang} files
  in the project folder are listed. &
  Write a \file{.disang} file, import it, and check with option~3. \\
 &
  The restraint mask helper counts ions as residues. &
  Use the protein's own residue numbers. \\
 &
  Cluster (SLURM) scripts may request a GPU for \texttt{sander}, use the wrong
  topology for constant pH, or contain an unfilled job dependency. &
  Read and edit the scripts before submitting. \\
\midrule
\key{5b} Analysis &
  Any list of analyses containing the digit 0 means ``exit'', so Angle (10)
  and Density maps (20) cannot be run. &
  Leave them out. \\
 &
  An atom that does not exist (CB of glycine) ends the analysis with a
  technical report. An empty solute mask is accepted. &
  Restart with valid atoms and masks. \\
 &
  The default SASA region (C-alpha atoms) gives a meaningless value. A heating
  stage is detected as NVT. &
  Choose region~4. Correct the type at \prompt{Confirm simulation type}. \\
 &
  Exports overwrite the same file names. Plots appear only as text on the
  screen. &
  Copy files you want to keep. Plot from the CSV file. \\
\midrule
\key{5c} QM/MM Preparator &
  Site residues are matched to the topology by number, so the HIGH layer can
  include a water and miss a propionate (1HRC). &
  Check the residue list by name and correct it (Chapter~\ref{ch:qmmm}). \\
 &
  The multiplicity defaults to 1 and is not checked. ScaleCharge values are
  not checked, and an old ScaleCharge stays after a change. &
  Set the multiplicity (2 for low-spin iron(III)). Check the route line. \\
 &
  Declining \prompt{Resume from saved state?} still reuses the old settings.
  A charge that is not a whole number only gives a warning. &
  Use option~3 (reset) first. Keep the total charge whole. \\
 &
  The MEDIUM layer method is fixed (HF/3-21G). &
  Edit the \file{.com} file if needed. \\
 &
  Extracting several frames fails. The frame syntax shows a symbol instead of
  \texttt{:end:}. &
  Extract one frame at a time, in a fresh session. \\
\midrule
\key{5d} Molecular Docking &
  A chain end without OXT is set up as charged, and preparing the receptor
  fails. &
  Set such ends to chain break (option~4). \\
 &
  A box around a ligand built from SMILES lies at the origin; overlap with the
  protein is not checked. &
  Define the box with \key{r} or \key{c}. \\
 &
  A \file{.smi} file cannot be used as a single ligand. Changing the ligand
  clears the box and kept groups. &
  Use \file{.sdf}, \file{.mol} or \file{.mol2}; set the box again. \\
 &
  With ad4, the top-scoring pose was not the closest to the crystal pose. &
  Start with vina; compare the first few poses. \\
\midrule
\key{ua} Structure Viewer &
  The viewer stops when ProPrep exits, and may use another port. Structures
  from different alignment runs may not overlap. &
  Use the address printed. Select the target at the end of each alignment. \\
\end{longtable}

\section{Sessions, templates and the command line}\label{sec:knownissues-sessions}

\begin{longtable}{@{}>{\raggedright\arraybackslash}p{0.16\linewidth}>{\raggedright\arraybackslash}p{0.44\linewidth}>{\raggedright\arraybackslash}p{0.34\linewidth}@{}}
\toprule
\textbf{Tool} & \textbf{What happens} & \textbf{Workaround} \\
\midrule
\endhead
\bottomrule
\endlastfoot
Sessions &
  \textsf{Continue from end} also replays a recorded exit, so ProPrep quits. &
  Delete the last two answers in the Session Editor. \\
 &
  Session files are found only in the folder ProPrep starts in. The default
  in the session list is \key{0} (cancel). &
  Start inside the project folder; type the session number. \\
 &
  \typed{undo} then \key{c} on a menu question replays the old answer. &
  Use \key{r}, or the Session Editor. \\
 &
  The menu-mode question is not recorded. &
  Add \texttt{--full-menu} to replay commands. \\
 &
  No \prompt{Session saved to:} line at a normal exit. &
  The file is saved as you go. \\
\midrule
Templates and batch &
  Variables are not detected automatically. &
  Mark them with \typed{v} in the Template Creator. \\
 &
  \texttt{--output-dir} does not create the folder, and every run fails. &
  Create the folder first. \\
 &
  Recorded choices are not checked, so a wrong value can give a wrong result
  without an error. &
  Check list values by eye. \\
 &
  A plain list of UniProt accessions is read as having a header. &
  Put \typed{input\_protein} on the first line. \\
 &
  Failed runs print a technical report over the progress bar. Forcing
  \texttt{--web} on a Terminal recording can break the replay. &
  Read the run's \file{error.log}. Let ProPrep choose the mode. \\
\midrule
Command line &
  \texttt{--output-dir} and \texttt{--session-dir} with a relative path are
  resolved in the wrong place; a missing session folder stops ProPrep. &
  Give a full path to a folder that exists. \\
 &
  \texttt{--pdbview}, \texttt{--force-replay} and \texttt{--analysis} still ask
  the menu-mode question. &
  Add \texttt{--full-menu}. \\
 &
  With both \texttt{--pdbid} and \texttt{--pdbfile}, the file is ignored. A
  failed \texttt{--pdbid} download is followed by a ``Downloaded and loaded''
  line. &
  Give only one. Trust \prompt{Failed to download}. \\
\end{longtable}

\section{Reporting new problems}\label{sec:knownissues-report}

Later versions of ProPrep may fix many of these limitations. Type
\texttt{proprep --version} to see which version you have, and check the
release notes when you update.

If you meet a problem that is not listed here or in
Chapter~\ref{ch:troubleshooting}, please tell the developers. Chapter~\ref{ch:help}
explains how to send a report and what to put in it.
